\begin{lemma}[Integration by parts against the weight $e^{-\sigma^2}$ on the line]
Let $g, g' : \mathbb{R} \to \mathbb{R}$ be such that $g$ is differentiable at
every $\sigma \in \mathbb{R}$ with derivative $g'(\sigma)$.  Assume the three
integrability conditions
\begin{enumerate}
  \item $\sigma \mapsto e^{-\sigma^2}\, g'(\sigma)$ is integrable on
    $\mathbb{R}$;
  \item $\sigma \mapsto e^{-\sigma^2}\,\bigl(\sigma\, g(\sigma)\bigr)$ is
    integrable on $\mathbb{R}$;
  \item $\sigma \mapsto e^{-\sigma^2}\, g(\sigma)$ is integrable on
    $\mathbb{R}$.
\end{enumerate}
Then
\[
  \int_{\mathbb{R}} e^{-\sigma^2}\,\sigma\, g(\sigma) \, d\sigma
    = \frac{1}{2} \int_{\mathbb{R}} e^{-\sigma^2}\, g'(\sigma) \, d\sigma .
\]

Both integrals are ordinary Lebesgue integrals over $\mathbb{R}$ with respect to
$d\sigma$: the weight $e^{-\sigma^2}$ sits inside the integrand, and the measure
is Lebesgue measure, unnormalised.  In particular this is not an identity about
the Gaussian measure $\gamma_1$ of the surrounding development, which carries an
extra normalising factor $\pi^{-1/2}$ appearing nowhere above; no Gaussian
measure occurs in this statement at all.  Note also that no decay hypothesis on
$g$ is assumed beyond the three integrability conditions --- there is no
requirement that $g$ vanish at infinity, and no boundary term appears in the
conclusion.
\end{lemma}
\begin{proof}
Set
\[
  G(\sigma) = -\tfrac{1}{2}\, e^{-\sigma^2} ,
\]
which is the point of the argument: $G$ is an explicit antiderivative of the
function $\sigma \mapsto \sigma\, e^{-\sigma^2}$ appearing on the left.

Step 1 ($G' (\sigma)= \sigma e^{-\sigma^2}$).  The map $s \mapsto -s^2$ is
differentiable with derivative $-2\sigma$ at $\sigma$, so by the chain rule
$s \mapsto e^{-s^2}$ is differentiable with derivative $-2\sigma e^{-\sigma^2}$,
and multiplying by the constant $-\tfrac12$,
\[
  G'(\sigma) = -\tfrac12 \cdot \bigl(-2\sigma\, e^{-\sigma^2}\bigr)
    = \sigma\, e^{-\sigma^2} \qquad (\sigma \in \mathbb{R}).
\]
In particular $G$ is differentiable at every point, as is $g$ by hypothesis.

Step 2 (recasting the integrability hypotheses).  The integration-by-parts
theorem to be applied asks for the integrability of the three products
$g' \cdot G$, $g \cdot G'$ and $g \cdot G$.  Each is a constant multiple of one
of the given integrands, by Step 1 and the definition of $G$: $g' \cdot G$ is
$-\tfrac12$ times the integrand of (1), $g \cdot G'$ is the integrand of (2)
rearranged, and $g \cdot G$ is $-\tfrac12$ times the integrand of (3).  Since a
constant multiple of an integrable function is integrable, all three products
are integrable.

Step 3 (integration by parts).  Both $g$ and $G$ are differentiable at every
real number, so in particular each is differentiable at every point of the
closed support of the other, and the integration-by-parts identity on
$\mathbb{R}$ applies with the integrability of Step 2 as its hypothesis:
\[
  \int_{\mathbb{R}} g(\sigma)\, G'(\sigma)\, d\sigma
    = -\int_{\mathbb{R}} g'(\sigma)\, G(\sigma)\, d\sigma .
\]
It is the integrability assumptions, not any decay of $g$, that licenses this;
no boundary contribution survives.

Step 4 (reading off the two sides).  By Step 1 the left-hand integrand is
$g(\sigma)\,\sigma\, e^{-\sigma^2}$, which is the integrand of the asserted
left-hand side, so
\[
  \int_{\mathbb{R}} e^{-\sigma^2}\,\sigma\, g(\sigma)\, d\sigma
    = \int_{\mathbb{R}} g(\sigma)\, G'(\sigma)\, d\sigma .
\]
On the right, substituting $G(\sigma) = -\tfrac12 e^{-\sigma^2}$ and taking the
constant out of the integral,
\[
  -\int_{\mathbb{R}} g'(\sigma)\, G(\sigma)\, d\sigma
    = -\int_{\mathbb{R}} g'(\sigma)\cdot
        \bigl(-\tfrac12 e^{-\sigma^2}\bigr)\, d\sigma
    = \frac{1}{2}\int_{\mathbb{R}} e^{-\sigma^2}\, g'(\sigma)\, d\sigma .
\]
Chaining these two identities with Step 3 gives the assertion.

Remark.  The identity is the Gaussian integration-by-parts formula in disguise:
dividing both sides by $\sqrt{\pi}$ turns the weight $e^{-\sigma^2}\,d\sigma$
into the probability measure $\gamma_1$, and the statement becomes
$\int \sigma\, g \, d\gamma_1 = \tfrac12 \int g' \, d\gamma_1$.  That normalised
form is not what is proved here, and the factor $\tfrac12$ above reflects the
variance $\tfrac12$ of the weight rather than any choice of normalisation.
\end{proof}
